\section*{Chapitre \ref{ch:nw:private-connectivity-to-crypto-venues} --- Connectivité privée vers les plateformes crypto}

\begin{solution}{exo:nw:private-connectivity-to-crypto-venues:1}
$0.014 \times 730 = 10.22$~USD par mois.
\end{solution}

\begin{solution}{exo:nw:private-connectivity-to-crypto-venues:2}
L'appairage relie le réseau de la plateforme à celui de chaque client~: la plateforme doit coordonner les espaces d'adressage et les routes avec chacun d'eux et fait confiance à leur
routage. Un service à point de terminaison publie un seul répartiteur de charge pour un nombre quelconque de clients, qui reçoivent chacun une interface dans leur propre réseau~; le
réseau de la plateforme reste fermé.
\end{solution}

\begin{solution}{exo:nw:private-connectivity-to-crypto-venues:3}
Le \omterm{def:nw:trading-in-a-public-cloud:pg}{groupe de placement} décide où sont lancées les instances du teneur de marché~: dans le même segment à forte bissection que les machines de la plateforme. Un point de terminaison
ne décide que de la route à partir de l'endroit où se trouve déjà l'instance.
\end{solution}

\begin{solution}{exo:nw:private-connectivity-to-crypto-venues:4}
Le point de terminaison n'a pas d'interface dans votre zone~: chaque paquet franchira une frontière de zone. Il faut vérifier où se trouvent les passerelles de la plateforme (leur
identifiant de zone)~; si elles sont dans az1 ou az2, il faut y déplacer vos instances~; si la plateforme peut ajouter une interface dans az4 et que ses passerelles sont dans az4, il faut le demander.
\end{solution}

\begin{solution}{exo:nw:private-connectivity-to-crypto-venues:5}
Il quitte le réseau privé vers la périphérie du fournisseur, aboutit à un \omterm{def:nw:the-european-map:pop}{point de présence} et revient à l'origine en traversant davantage d'équipements
(et généralement une terminaison TLS et un répartiteur de charge)~; dans le modèle, cela coûte \qty{2}{\milli\second} contre quelques dizaines de microsecondes pour les
sauts privés.
\end{solution}

\begin{solution}{exo:nw:private-connectivity-to-crypto-venues:6}
Parce que, dans ce bâtiment, l'\omterm{def:nw:colocation-products-and-how-they-are-sold:mmr}{interconnexion} (cross-connect) n'est que quelques mètres de fibre menant droit au réseau de la plateforme, alors que PrivateLink passe
par la région du fournisseur de cloud et revient~; la documentation propre à la plateforme recommande les \omterm{def:nw:colocation-products-and-how-they-are-sold:mmr}{interconnexions} pour la latence la plus faible et la plus stable.
\end{solution}

\begin{solution}{exo:nw:private-connectivity-to-crypto-venues:7}
Environ \qty{126}{\micro\second} de sauts au point de terminaison~: la médiane du point de terminaison de même zone atteint alors le 99\up{e}~centile de l'appairage simple, soit
\qty{396}{\micro\second}.
\end{solution}

\begin{solution}{exo:nw:private-connectivity-to-crypto-venues:8}
Un \omterm{def:nw:private-connectivity-to-crypto-venues:endpoint}{point de terminaison privé} est privé, pas forcément rapide~: il ajoute des sauts, il peut être servi depuis une autre zone, et un \omterm{def:nw:trading-in-a-public-cloud:pg}{groupe de placement} partagé ou un simple appairage,
lorsque la plateforme les propose, sont plus rapides. «~Privé~» indique qui peut voir le trafic, non combien de temps il met.
\end{solution}

\begin{solution}{pb:nw:private-connectivity-to-crypto-venues:1}
\textbf{Partie I.}
\begin{enumerate}
\item 310 et \qty{436}{\micro\second}.
\item 595 et \qty{797}{\micro\second}.
\item \qty{285}{\micro\second} jusqu'à la médiane, \qty{360}{\micro\second} jusqu'au 99\up{e}~centile.
\item Sept fois l'écart de \qty{40}{\micro\second} entre l'appairage et un point de terminaison de même zone.
\end{enumerate}
\textbf{Partie II.}
\begin{enumerate}[resume]
\item 0.004~USD par le point de terminaison de même zone, 0.012 par l'autre.
\item Tout le transfert inter-zones est facturé~: $0.4~\mathrm{GB} \times 0.01 \times 2$ sens $= 0.008$~USD.
\item $10.22 + 40 = 50.22$~USD et $10.22 + 120 = 130.22$~USD par mois.
\item 10.22~USD dans une seule zone, 30.66 dans trois.
\end{enumerate}
\textbf{Partie III.}
\begin{enumerate}[resume]
\item En comparant les identifiants de zone de ses instances à ceux des interfaces du point de terminaison, et en mesurant~: une médiane voisine de \qty{600}{\micro\second}
  là où l'on attendait 300 est le signe.
\item Déplacer ses instances dans une zone où le point de terminaison a une interface (si les passerelles de la plateforme s'y trouvent), ou demander à la plateforme une interface dans
  sa zone, ou emprunter un autre chemin proposé par la plateforme (appairage, \omterm{def:nw:trading-in-a-public-cloud:pg}{groupe de placement} partagé).
\item Le segment propre à la plateforme dans le réseau~: environ \qty{150}{\micro\second} d'aller-retour médian dans le modèle, le meilleur de tous.
\item Le déplacement du point de terminaison n'aide que si les passerelles de la plateforme sont dans la zone de la firme~; sinon la frontière est franchie quand même, et les instances
  doivent aussi être déplacées.
\end{enumerate}
\textbf{Partie IV.}
\begin{enumerate}[resume]
\item \emph{Résultat nommé}~: servi depuis une autre zone, le \omterm{def:nw:private-connectivity-to-crypto-venues:endpoint}{point de terminaison privé} ajoute environ \qty{285}{\micro\second} à l'aller-retour médian (et
  \qty{360}{\micro\second} au 99\up{e}~centile) et 0.008~USD par million de messages de 400~octets en transfert inter-zones, en plus de 0.004~USD de
  traitement des données.
\item Documentation~: les chemins, le comportement par zone, les tarifs. Mesure~: les allers-retours du chapitre~16. Hypothèses~: les sauts du point de terminaison, la taille des messages
  et la médiane du cluster partagé.
\item Parce qu'il coûte quelques centimes par million de messages et des dizaines de dollars par mois~: son prix est négligeable au regard de ce que valent les microsecondes.
\item Les identifiants de zone de ses passerelles et des interfaces de son point de terminaison, par environnement.
\item Après chaque relance, après la maintenance de la plateforme, et en continu par la mesure (chapitre~17).
\item Les deux cas portent sur l'endroit où passe le chemin du client par rapport à la plateforme~: dans une salle, la plateforme égalise les chemins~; dans une région cloud, personne ne le fait, et
  le client doit trouver le plus court.
\item Le chemin serait une \omterm{def:nw:colocation-products-and-how-they-are-sold:mmr}{interconnexion} dans le bâtiment (chapitre~9), et la question de la zone disparaîtrait.
\item Il garantit que le trafic reste sur le réseau privé du fournisseur, et non qu'il prenne le chemin le plus court ou qu'il reste dans une seule zone.
\end{enumerate}
\end{solution}

\begin{solution}{iq:nw:private-connectivity-to-crypto-venues:1}
L'appairage relie deux réseaux entiers~; PrivateLink expose un seul service (derrière un répartiteur de charge) au moyen d'une interface dans le réseau du client,
sans routage entre eux. L'appairage a moins de sauts~; PrivateLink passe à l'échelle pour de nombreux clients et garde fermé le réseau du fournisseur.
\iqlookfor{Routage contre exposition de service~; sauts~; qui contrôle quoi.}
\end{solution}

\begin{solution}{iq:nw:private-connectivity-to-crypto-venues:2}
On compare les identifiants de zone (et non les noms) des instances à ceux des interfaces du point de terminaison et des passerelles de la plateforme~; on vérifie le réglage
inter-zones du répartiteur de charge~; et on mesure~: un aller-retour inter-zones est visiblement plus long.
\iqlookfor{Identifiants~; comportement du répartiteur de charge~; mesure.}
\end{solution}

\begin{solution}{iq:nw:private-connectivity-to-crypto-venues:3}
Selon le besoin de latence, face au coût et à l'éligibilité~: un \omterm{def:nw:trading-in-a-public-cloud:pg}{groupe de placement} partagé ou l'appairage si la stratégie joue sur la vitesse, un point de terminaison si elle a besoin de confidentialité et
d'une vitesse raisonnable, le chemin public sinon. Il faut demander les zones, les limites de débit par palier, la capacité des passerelles et les conditions d'éligibilité.
\iqlookfor{Paliers associés aux stratégies~; zones~; limites de débit~; conditions.}
\end{solution}

\begin{solution}{iq:nw:private-connectivity-to-crypto-venues:4}
La pile de l'instance, l'interface du point de terminaison, le répartiteur de charge, tout saut inter-zones, la passerelle de la plateforme. On horodate matériellement les paquets sur
l'instance, on les compare aux horodatages d'accusé de réception de la plateforme (avec précaution sur son horloge, chapitre~20), et on fait varier un élément à la fois (zone, chemin).
\iqlookfor{Décomposition~; horodatages matériels~; expériences contrôlées.}
\end{solution}

\begin{solution}{iq:nw:private-connectivity-to-crypto-venues:5}
Elle est équitable si elle est proposée à tout participant qui remplit les conditions, selon des termes publiés et non discriminatoires, comme doit l'être la \omterm{def:nw:colocation-products-and-how-they-are-sold:colo}{colocation} (chapitre~9)~; elle ne l'est pas
si elle relève d'un arrangement privé. La question porte sur les conditions de la plateforme, non sur la technologie.
\iqlookfor{Conditions égales au sein d'un même service~; termes publiés~; l'analogie avec la \omterm{def:nw:colocation-products-and-how-they-are-sold:colo}{colocation}.}
\end{solution}

\begin{solution}{iq:nw:private-connectivity-to-crypto-venues:6}
Des instances dans les deux zones, chacune connectée aux passerelles de la plateforme dans sa propre zone (interfaces de point de terminaison ou appairage dans chacune), avec un routage des ordres qui
privilégie la zone du moteur d'appariement principal de la plateforme (matching engine) et bascule sur l'autre~; on mesure les deux chemins en continu et on teste le basculement.
\iqlookfor{Alignement par zone~; chemins actifs dans les deux zones~; basculement~; mesure.}
\end{solution}
